\section{Background}\label{sec:background}

\paragraph{Verification, identification, and Eigenfaces.} \emph{Identification} asks which enrolled identity a probe belongs to; \emph{verification} asks whether it matches a claimed identity. Both need a representation $\phi$ and a distance $d$; verification compares $d(\phi(x),\phi(y))$ with a threshold $\tau$. In the open-set setting the claimed identity may not be enrolled, so the system must say ``unknown''---the setting of face unlock. Our verification experiment is closed-set: all 40 identities are enrolled, so it measures cross-identity confusion rather than rejection of unknown people. Eigenfaces stacks each centred image as a vector in $\mathbb{R}^{4096}$, computes the leading eigenvectors of the covariance, and projects onto the first $k$ \citep{turk1991}. The components can be reshaped into images, but the objective is variance, not between-identity separation, and on faces the largest variance is usually illumination. On tightly cropped, roughly aligned data such as Olivetti it is nevertheless a strong baseline, and therefore the right comparator.

\paragraph{Deep face representations.} DeepFace showed that a deep network trained on millions of faces could approach human verification performance \citep{taigman2014}. FaceNet's triplet loss optimises the verification property directly: anchor--positive close, anchor--negative far by a margin \citep{schroff2015}. VGGFace2 scaled the labelled corpus to 3.31 million images of 9,131 identities \citep{cao2018}; ArcFace replaced the triplet objective with an additive angular margin \citep{deng2019}; surveys trace the progression \citep{wang2021}. Here the network is a fixed feature extractor: Inception-ResNet-v1 \citep{szegedy2017} pretrained on VGGFace2 via \texttt{facenet-pytorch} \citep{facenetpytorch}, frozen, producing an L2-normalised 512-d vector. Freezing isolates what a pretrained representation contains without confounding the comparison with training that 400 images could not support. For unit vectors, cosine similarity is $1-d_{\mathrm{eucl}}^2/2$.

\paragraph{Detection and illumination.} MTCNN is a cascade: a proposal network (P-Net) scans an image pyramid, a refinement network (R-Net) filters and regresses, and an output network (O-Net) refines and predicts five landmarks \citep{zhang2016}. The \texttt{facenet-pytorch} implementation crops the O-Net box to a square and resizes to $160\times160$; it returns the landmarks but does not warp the patch, so alignment is detector-box alignment. Olivetti avoids the issue: its faces are pre-aligned, which makes it the best case for raw pixels and an easy case for the deep model. Illumination is the canonical nuisance: Extended Yale Face Database B held pose and expression roughly constant while sweeping light, and illumination-cone models showed that images under all lighting directions lie near a low-dimensional convex cone \citep{georghiades2001}. Changing the \emph{amount} of light scales pixel values roughly uniformly; changing the \emph{direction} produces shading that moves across the face.

\paragraph{Metrics and visualisation.} Verification is characterised by false accept rate (FAR), the fraction of impostor trials accepted, and false reject rate (FRR), the fraction of genuine trials rejected. Sweeping the threshold traces an ROC curve of TAR $=1-\mathrm{FRR}$ against FAR; summary points are the equal error rate (EER), where FAR $=$ FRR (interpolated here from a finite score grid), and TAR at fixed FAR, which is how operational systems are compared because deployments fix the acceptable false-accept rate rather than the threshold \citep{grother2019}. FERET established the vocabulary of defined gallery/probe splits and impostor trials \citep{phillips2000}. We also use discriminability $d'$ under an equal-variance Gaussian assumption,
\begin{equation}
d'=\frac{\mu_{\mathrm{imp}}-\mu_{\mathrm{gen}}}{\sqrt{(\sigma_{\mathrm{gen}}^2+\sigma_{\mathrm{imp}}^2)/2}}.
\label{eq:dprime}
\end{equation}
PCA is a projection, and a projection can lie by omission; t-SNE is the usual cautionary example \citep{maaten2008}, but ordinary PCA scatters have a subtler problem: the two directions of largest variance need not be the two that separate identities. We quantify how much of the full-dimensional geometry a 2-D map preserves using the Spearman rank correlation between 2-D and full-dimensional pairwise distances, and never compute a metric on 2-D coordinates.
